% M7 semantic reconstruction of physical page 114 / printed page 113.
\begin{HLAClauseSemanticPage}{113}
\HLAStableID{m7.page.114}
\HLAStableID{m7.p114.b01}
\HLAClausePlacedBlock{90.000}{68.648}{56.728}{%
  \HLAAnnexHeading{Annex D}
}
\HLAStableID{m7.p114.b02}
\HLAClausePlacedBlock{90.000}{97.140}{63.882}{%
  \HLAAnnexStatus{(normative)}
}
\HLAStableID{m7.p114.b03}
\HLAClausePlacedBlock{90.000}{128.648}{194.403}{%
  \HLAAnnexTitle{OMT data interchange format}
}
\HLAStableID{m7.p114.b04}
\HLAClausePlacedBlock{90.000}{157.806}{432.087}{%
  \HLAClauseParagraph{\resizebox{432.075bp}{!}{\mbox{The HLA OMT data interchange format (DIF) is a standard data exchange format used to store and transfer}}\newline{}\resizebox{432.047bp}{!}{\mbox{FOMs, SOMs, and HLA object model components (e.g., the HLA object modeling aspects of BOMs, as well}}\newline{}\resizebox{431.945bp}{!}{\mbox{as FOM modules, SOM modules, and MIMs) between object modeling tools and to initialize the HLA RTI.}}\newline{}\resizebox{432.087bp}{!}{\mbox{The DIF is built on a common metamodel that represents the information needed to represent and manage}}\newline{}\resizebox{212.599bp}{!}{\mbox{object models created using the HLA OMT standard.}}}
}
\HLAStableID{m7.p114.b05}
\HLAClausePlacedBlock{90.000}{241.140}{129.254}{%
  \HLAAnnexSubclauseHeading{D.1 Notation of the DIF}
}
\HLAStableID{m7.p114.b06}
\HLAClausePlacedBlock{90.000}{267.786}{432.000}{%
  \HLAClauseParagraph{\resizebox{432.000bp}{!}{\mbox{In order to ensure that the definition of the DIF is unambiguous, it has been formally defined in terms of}}\newline{}\resizebox{431.959bp}{!}{\mbox{XML, a formal notation used to describe document specifications. The standard for XML notation is}}\newline{}\resizebox{141.786bp}{!}{\mbox{described in REC-XML-20060816.}}}
}
\HLAStableID{m7.p114.b07}
\HLAClausePlacedBlock{90.000}{327.120}{216.408}{%
  \HLAAnnexSubclauseHeading{D.2 OMT DIF XML Schema declaration}
}
\HLAStableID{m7.p114.b08}
\HLAClausePlacedBlock{90.000}{367.806}{432.000}{%
  \HLAClauseParagraph{\resizebox{431.980bp}{!}{\mbox{XML calls for using an XML Schema or a Document Type Declaration (DTD) to specify the valid syntax,}}\newline{}\resizebox{431.967bp}{!}{\mbox{structure, and format for documents. The XML Schema that describes the OMT DIF is included in the}}\newline{}\resizebox{92.431bp}{!}{\mbox{paragraphs that follow.}}}
}
\HLAStableID{m7.p114.b09}
\HLAClausePlacedBlock{90.000}{415.805}{432.039}{%
  \HLAClauseParagraph{\resizebox{431.999bp}{!}{\mbox{The XML specification defines two levels of correctness for XML documents: “well-formed” and “valid.” A}}\newline{}\resizebox{432.018bp}{!}{\mbox{well-formed XML document is one that follows the syntax constraints of XML; for example, all element}}\newline{}\resizebox{431.938bp}{!}{\mbox{opening and closing tags are appropriately nested. A valid XML file is one that is well formed and complies}}\newline{}\resizebox{431.980bp}{!}{\mbox{with an XML Schema. That is, a valid XML document shall contain all elements and attributes defined in}}\newline{}\resizebox{432.038bp}{!}{\mbox{the XML Schema and only those elements and attributes. To support the exchange of object models among}}\newline{}\resizebox{351.428bp}{!}{\mbox{heterogeneous tools, all OMT DIF-compliant files shall be valid, and thus, well formed.}}}
}
\HLAStableID{m7.p114.b10}
\HLAClausePlacedBlock{90.001}{499.804}{432.047}{%
  \HLAClauseParagraph{\resizebox{431.969bp}{!}{\mbox{Because it is often necessary to exchange incomplete OMT documents using the OMT DIF (for federations}}\newline{}\resizebox{432.007bp}{!}{\mbox{under development, for example), many completeness and consistency constraints are not included in the}}\newline{}\resizebox{431.823bp}{!}{\mbox{OMT DIF XML Schema. Therefore, much of the individual information is specified as optional in the XML}}\newline{}\resizebox{432.047bp}{!}{\mbox{Schema. However, when complete, an OMT document, in addition to being a valid XML document, shall be}}\newline{}\resizebox{432.021bp}{!}{\mbox{consistent and complete with respect to the body of this document. In the OMT XML Schema, note labels}}\newline{}\resizebox{431.988bp}{!}{\mbox{and note references shall be prefixed with an underscore to ensure that the labeling conforms to the}}\newline{}\resizebox{432.000bp}{!}{\mbox{NCName data type convention. Also, several instances exist in which the tabular representation differs from}}\newline{}\resizebox{144.629bp}{!}{\mbox{the XML representation, as follows:}}}
}
\HLAStableID{m7.p114.b11}
\HLAClausePlacedBlock{100.022}{601.802}{424.078}{%
  \HLAClauseDashItem{\resizebox{400.006bp}{!}{\mbox{In the tabular representation of object models, object class and interaction class inheritance are}}\newline{}\resizebox{400.042bp}{!}{\mbox{represented by columns in the which the parent of a class was represented to its left and children to}}\newline{}\resizebox{400.114bp}{!}{\mbox{its right. In the XML depiction of object models, inheritance is represented by membership, in which}}\newline{}\resizebox{392.069bp}{!}{\mbox{children of a parent class are depicted as class elements contained within the parent class element.}}}
}
\HLAStableID{m7.p114.b12}
\HLAClausePlacedBlock{100.021}{652.801}{424.079}{%
  \HLAClauseDashItem{\resizebox{400.117bp}{!}{\mbox{In the tabular representation, multiple entries in a field are represented by comma-separated lists. In}}\newline{}\resizebox{400.042bp}{!}{\mbox{XML, elements are separated by spaces. This situation exists in notes and attributes’ available}}\newline{}\resizebox{48.027bp}{!}{\mbox{dimensions.}}}
}
\HLAStableID{m7.p114.b13}
\HLAClausePlacedBlock{100.021}{691.800}{424.079}{%
  \HLAClauseDashItem{\resizebox{400.067bp}{!}{\mbox{The formatting of notes in XML is different from that presented in the tabular representation in the}}\newline{}\resizebox{400.029bp}{!}{\mbox{body of this specification. In the tabular representation, a pointer to a note is indicated by an}}}
}
\end{HLAClauseSemanticPage}
